\documentclass[10pt,twocolumn,letterpaper]{article}

% Official ACCV 2026 rebuttal formatting. Keep these settings unchanged.
\usepackage[rebuttal]{cvpr}
\usepackage{accvabbrv}
\usepackage{graphicx}
\usepackage{booktabs}
% Optional in the official template; incompatible with the local Tectonic engine.
% Re-enable when compiling with the conference-recommended pdflatex toolchain.
% \usepackage[accsupp]{axessibility}
\definecolor{accvblue}{rgb}{0.12,0.49,0.85}
\usepackage[pagebackref,breaklinks,colorlinks,citecolor=accvblue]{hyperref}

\def\paperID{1107}
\def\confName{ACCV}
\def\confYear{2026}

\begin{document}
\title{Exploring RGB and Video V-JEPA 2.1 Features for Aerial Depth Estimation and Semantic Scene Completion}
\maketitle
\thispagestyle{empty}
\appendix

% Response strategy: (1) repair trust/presentation, (2) state the exact positive result,
% (3) narrow unsupported causal claims, and (4) explain the negative SSC/video findings.
\noindent\textbf{We thank LD5V, CWVJ, and y7qY for their detailed feedback.}
We agree that unresolved appendix references and several presentation errors obscured the
submitted evidence. We clarify the main result and appropriately narrow our conclusions
below; no unrequested new result is introduced.

\paragraph{Presentation and protocol corrections (all reviewers).}
The appendix tables/figures were not packaged with resolvable references; we apologize
because this made supporting evidence inaccessible. Table~3's caption incorrectly names
baselines absent from its body and highlights only RMSE; it should name only the included
DepthAnything2-OccuFly baseline and compare every metric. Table~4 is a stress-test result,
not a comparative ablation. Equation~19 is correct: training uses 30/50\,m, while 40\,m
validation uses scenes 01--04 and testing uses 05--08; ``1--9'' in Table~4 is an error.
Because these scene identities occur at other training altitudes, this holds out altitude,
not scene identity. We therefore retract the scene-generalization interpretation. We will
repair all references, cite every figure, verify figure text size, and replace the
screenshot-like overview with an end-to-end data-flow diagram distinguishing frozen and
trainable components.

\paragraph{What the depth comparison establishes (CWVJ, y7qY).}
FiLM-DPT-RGB versus DepthAnything2-OccuFly improves RMSE (4.621 vs.~4.844\,m), MAE
(3.523 vs.~4.193\,m), and AbsRel (0.129 vs.~0.134), is essentially tied but slightly worse
on $\delta_1$ (0.833 vs.~0.834), and is worse on SILog (0.164 vs.~0.112). RMSE/MAE measure
absolute metric error (with RMSE emphasizing large errors), whereas $\delta$ and SILog
measure multiplicative/log-space consistency. Thus the supported conclusion is
\emph{competitive absolute metric accuracy from a frozen general-purpose backbone}, not
overall superiority. The metric pattern is an observation; without error-distribution
analysis we do not assert its mechanism.

\paragraph{FiLM, backbone, and uncertainty scope (LD5V, CWVJ).}
We agree that the unexpectedly weak plain-DPT row prevents a clean causal attribution to
FiLM. We consequently restrict the claim to the \emph{configured FiLM-DPT system} being
the strongest tested V-JEPA head; Table~4 alone does not isolate conditioning. Likewise,
one frozen encoder cannot demonstrate V-JEPA-specific superiority, so conclusions are
limited to V-JEPA~2.1 under this OccuFly protocol. Matched DPT/FiLM conditioning ablations,
a second frozen backbone, and multi-seed uncertainty are the correct requested controls;
we do not present unavailable results as established evidence.

\paragraph{Video, MVS, and SSC interpretation (LD5V, CWVJ, y7qY).}
The strongest submitted SSC result is 37.68 SC IoU / 3.46 SSC mIoU; the strongest
Pose-MVS result is 32.55 / 2.36. Hence video/MVS does not improve the simpler RGB-context
configuration, and we retract any implication of a new superior SSC architecture. DTA
does not explicitly warp temporal features into the target camera frame, so UAV ego-motion
can cause spatial mixing. Pose-MVS was evaluated only through downstream SSC, not direct
depth; it is therefore an exploratory negative result, not evidence that MVS depth itself
fails. Direct depth, pose-aligned temporal aggregation, oracle lifting, per-class analysis,
and repeated seeds are needed to localize the bottleneck. Finally, eight H100s describe
experimental infrastructure, not demonstrated onboard deployment; we will remove that
implication and report efficiency before making it.

\end{document}
